% =============================================================================
% TEMPLATE DISPENSE UNIVERSITARIE - CLASSE REPORT (GENERIC & READY-TO-USE)
% =============================================================================
% ISTRUZIONI COMPILAZIONE (pdflatex):
% 1. Compilare due volte per risolvere riferimenti incrociati e indici
% 2. Creare cartella 'assets/' nella stessa directory del file .tex
% 3. Per indice analitico: eseguire 'makeindex main.tex' se si usa \index{}
% =============================================================================
\documentclass[11pt,a4paper,italian,twoside]{report}

% ==========================================
% 1. CODIFICA, LINGUA E PRECARICAMENTO SICURO
% ==========================================
\usepackage[utf8]{inputenc}
\usepackage[T1]{fontenc}
\usepackage[italian]{babel}
\usepackage{microtype} % Ottimizzazione tipografica avanzata

% ==========================================
% 2. PACCHETTI MATEMATICI (PRIMA DEI FONT)
% ==========================================
\usepackage{amsmath,amsthm,amssymb,amsfonts}
\usepackage{mathtools,bm} % mathtools estende amsmath, bm gestisce grassetto vettoriale

% ==========================================
% 3. FONT (DOPO AMS: PREVIENE CONFLITTI INTERNI)
% ==========================================
\usepackage{newtxtext}
\usepackage[amsthm]{newtxmath} % [amsthm] essenziale per compatibilità

% ==========================================
% 4. IMPAGINAZIONE E INTERLINEA
% ==========================================
\usepackage[
  inner=25mm, outer=20mm,
  top=20mm, bottom=20mm,
  headheight=15pt, footskip=15pt
]{geometry}
\usepackage{setspace}
\setstretch{1.15} % Interlinea ottimizzata per lettura prolungata


% ==========================================
% RIDUZIONE SPAZIO VUOTO PRIMA DEI CAPITOLI
% ==========================================
\usepackage{etoolbox}
\makeatletter
\patchcmd{\@makechapterhead}{\vspace*{0\p@}}{\vspace*{0pt}}{}{} % Spazio prima del numero
\patchcmd{\@makechapterhead}{\vskip 20\p@}{\vskip 5pt}{}{}         % Spazio tra numero e titolo
\patchcmd{\@makechapterhead}{\vskip 40\p@}{\vskip 15pt}{}{}        % Spazio dopo il titolo
\makeatother
% ==========================================
% 5. COLORI E VARIABILI GLOBALI PERSONALIZZABILI
% ==========================================
\usepackage{xcolor}
\definecolor{primary}{HTML}{2C3E50}    % Colore primario (blu scuro)
\definecolor{secondary}{HTML}{34495E}  % Colore secondario (grigio-blu)
\definecolor{accent}{HTML}{E67E22}     % Colore accento (arancione sobrio)
\definecolor{code-bg}{HTML}{F4F6F7}    % Sfondo blocchi codice
\definecolor{text-main}{HTML}{212529}  % Colore testo principale

% MODIFICARE QUI I METADATI DEL CORSO
\newcommand{\CorsoTitolo}{Fisiologia dell'esercizio fisico}

% ==========================================
% 6. AMBIENTI TEOREMATICI E DEFINIZIONI
% ==========================================
% Numerazione automatica: [section] fa ripartire da 1 ad ogni nuova sezione
\newtheorem{theorem}{Teorema}[section]
\newtheorem{lemma}[theorem]{Lemma}
\newtheorem{corollary}[theorem]{Corollario}
\newtheorem{proposition}[theorem]{Proposizione}

\theoremstyle{definition} % Stile "normale" (non corsivo)
\newtheorem{definition}{Definizione}[section]
\newtheorem{example}{Esempio}[section]
\newtheorem{remark}{Osservazione}[section]
\newtheorem{note}{Nota}[section]

\renewcommand{\proofname}{Dimostrazione}

% ==========================================
% 7. TABELLE, FIGURE E FLOAT
% ==========================================
\usepackage{booktabs,tabularx,longtable,makecell,multirow} % aggiunto multirow
\usepackage{graphicx}
\graphicspath{{assets/}} % Percorso predefinito per le immagini
\usepackage{caption,subcaption,float}
\captionsetup{labelfont=bf,textfont=it,labelsep=period,font=small}
\restylefloat{table}
\restylefloat{figure}

% ==========================================
% 8. LISTE ED ENUMERAZIONI
% ==========================================
\usepackage{enumitem}
\setlist[itemize]{label=\textbullet,leftmargin=*,topsep=0.2\baselineskip,itemsep=0.1\baselineskip}
\setlist[enumerate]{label=\arabic*.,leftmargin=*,topsep=0.2\baselineskip,itemsep=0.1\baselineskip}

% ==========================================
% 9. CODICE SORGENTE (Listings)
% ==========================================
\usepackage{listings}
\lstset{
  basicstyle=\ttfamily\footnotesize,
  breaklines=true,
  backgroundcolor=\color{code-bg},
  keywordstyle=\color{primary}\bfseries,
  commentstyle=\itshape\color{gray},
  stringstyle=\color{secondary},
  numbers=left,
  numberstyle=\tiny\color{gray},
  tabsize=4,
  captionpos=b,
  xleftmargin=10pt,
  rulecolor=\color{secondary},
  framerule=0.4pt
}

% ==========================================
% 10. BOX PERSONALIZZATI (tcolorbox)
% ==========================================
% Ambienti già pronti: definizioni, teoremi, concetti, esempi, riassunti, avvisi
\usepackage{tcolorbox}
\tcbuselibrary{skins,breakable}

\tcbset{
  dispensebox/.style={
    breakable, enhanced, arc=2.5mm, boxrule=0.4pt,
    colframe=#1!75!black, colback=#1!4,
    left=5mm, right=5mm, top=5mm, bottom=5mm,
    attach boxed title to top left={xshift=6mm, yshift=-2mm},
    boxed title style={size=small, colback=#1!85!black, colupper=white, arc=1.5mm, boxrule=0pt, fonttitle=\bfseries}
  }
}

\newtcolorbox{definizione}{dispensebox=secondary,title=Definizione}
\newtcolorbox{teorema}{dispensebox=primary,title=Teorema}
\newtcolorbox{concetto}{dispensebox=primary!80!black,title=Concetto Importante}
\newtcolorbox{esempio}{dispensebox=green!60!gray,title=Esempio}
\newtcolorbox{riassunto}{dispensebox=accent!80!yellow,title=Riassunto}
\newtcolorbox{avviso}{dispensebox=red!80!black,title=Avviso}

% ==========================================
% 11. HEADER, FOOTER E NAVIGAZIONE
% ==========================================
\usepackage{fancyhdr}
\pagestyle{fancy}
\fancyhf{}
\fancyhead[LE,RO]{\thepage}
\fancyhead[RE]{\small\scshape \CorsoTitolo}
\fancyhead[LO]{\leftmark} % Mostra nome capitolo/section
\renewcommand{\chaptermark}[1]{\markboth{#1}{}}
\renewcommand{\sectionmark}[1]{\markright{#1}}
\renewcommand{\headrulewidth}{0.4pt}

\fancypagestyle{plain}{
  \fancyhf{}
  \fancyfoot[C]{\thepage}
  \renewcommand{\headrulewidth}{0pt}
}

% ==========================================
% 12. HYPERLINK, METADATA E CROSS-REF
% ==========================================
\usepackage{hyperref}
\usepackage{bookmark} % Gestione avanzata segnalibri PDF

% Previene warning hyperref per comandi matematici/formatting nei segnalibri PDF
\pdfstringdefDisableCommands{
  \let\textsubscript\relax
  \let\subscript\relax
  \def\textsuperscript#1{^#1}
  \def\textsubscript#1{_#1}
  \let\_\textunderscore
}

\hypersetup{
  colorlinks=true,
  linkcolor=primary,
  urlcolor=secondary,
  citecolor=secondary,
  pdfborder={0 0 0},
  pdftitle=\CorsoTitolo,
  pdfsubject={Dispense Universitarie},
  pdfkeywords={Università, Dispense, LaTeX}
}

\usepackage{cleveref} % Riferimenti automatici intelligenti (es. \cref{eq:1})
\usepackage{imakeidx}
\makeindex % Abilita indice analitico con \index{termine}

% ==========================================
% ALLINEAMENTO SPAZIATURA CAPITOLI
% ==========================================
\usepackage{etoolbox}
\makeatletter
% Rimuove lo spazio extra sopra il capitolo: allinea il titolo al margine superiore standard
\patchcmd{\@makechapterhead}{\vspace*{50\p@}}{\vspace*{0pt}}{}{}
% Riduce la distanza tra numero e titolo
\patchcmd{\@makechapterhead}{\vskip 20\p@}{\vskip 5pt}{}{}
% Riduce la distanza dopo il titolo (coerente con le sezioni successive)
\patchcmd{\@makechapterhead}{\vskip 40\p@}{\vskip 10pt}{}{}

% Applica la stessa modifica anche ai capitoli non numerati (\chapter*)
\patchcmd{\@makeschapterhead}{\vspace*{50\p@}}{\vspace*{0pt}}{}{}
\makeatother

% =============================================================================
% INIZIO DOCUMENTO
% =============================================================================
\begin{document}

% ==========================================
% TITLE PAGE
% ==========================================
\begin{titlepage}
  \centering
  \vspace*{2.5cm}
  {\Large\scshape \CorsoTitolo\par}
  \vspace{2.2cm}
\end{titlepage}

% ==========================================
% INDICE DEI CONTENUTI
% ==========================================
\tableofcontents

% ==========================================
% CONTENUTO PRINCIPALE
% ==========================================
% STRUTTURA CONSIGLIATA:
% \chapter{Titolo Capitolo}
% \section{Titolo Sezione}
% \subsection{Titolo Sottosezione}
%
% NOTA IMPORTANTE PER I SEGNALIBRI PDF:
% Se usi pedici/apici matematici nei titoli (\chapter, \section, ecc.), 
% hyperref potrebbe generare warning. Usa sempre \texorpdfstring{math}{text}:
% Esempio corretto: \chapter{Studio della molecola di \texorpdfstring{$H_2O$}{H2O}}
%
% PER INSERIRE IMMAGINI:
% \begin{figure}[htbp]
%   \centering
%   \includegraphics[width=0.8\textwidth]{nome_immagine.pdf}
%   \caption{Didascalia immagine}
%   \label{fig:etichetta}
% \end{figure}
%
% PER INSERIRE FORMULE:
% Ambiente align per allineare più equazioni
% \begin{align}
%   f(x) &= ax^2 + bx + c \\
%   f'(x) &= 2ax + b
% \end{align}
%
% PER INSERIRE TABELLE:
% \begin{table}[htbp]
%   \centering
%   \begin{tabularx}{\textwidth}{l X r}
%     \toprule
%     Colonna 1 & Contenuto flessibile & Allineata a destra \\
%     \midrule
%     Dato A & Testo lungo... & 12.5 \\
%     \bottomrule
%   \end{tabularx}
%   \caption{Didascalia tabella}
% \end{table}
%
% PER USARE I BOX PERSONALIZZATI:
% \begin{definizione} Testo della definizione... \end{definizione}
% \begin{teorema} Testo del teorema... \end{teorema}
% \begin{esempio} Testo dell'esempio... \end{esempio}
% \begin{proof} Passaggi dimostrativi... \end{proof}
%
% PER INDICIZZAZIONE (Indice Analitico):
% Scrivere \index{termine} nel punto desiderato del testo.
% =============================================================================

% Inserire qui i capitoli del corso.
% \chapter{Fondamenti Teorici}
% \section{Definizioni di base}
% Testo introduttivo...
% \begin{definizione}
% Definizione del concetto chiave...
% \end{definizione}

% =============================================================================
% CAPITOLO: OMEOSTASI
% Corso: Fisiologia dell'esercizio fisico
% =============================================================================

\chapter{Omeostasi}
\label{chap:omeostasi}

% -----------------------------------------------------------------------------
% 1. DEFINIZIONE E CONCETTI FONDAMENTALI
% -----------------------------------------------------------------------------
\section{Definizione di omeostasi}
\label{sec:definizione}

\begin{definizione}
\textbf{Omeostasi}: mantenimento di condizioni stabili dell'organismo (equilibrio corporeo, stato di benessere) mediante un insieme di processi e meccanismi fisiologici che correggono tempestivamente le perturbazioni del sistema.
\end{definizione}

È fondamentale distinguere:
\begin{itemize}
  \item \textbf{Ambiente esterno}: l'ambiente in cui vive l'organismo;
  \item \textbf{Ambiente interno}: l'ambiente in cui vivono gli elementi che formano l'organismo (plasma e liquidi corporei).
\end{itemize}

\begin{concetto}
Le condizioni omeostatiche non sono fisse: si tratta di un \textbf{equilibrio dinamico}, basato su stabilità, bilanciamento e adattamento continuo, non su immobilità.
\end{concetto}

% -----------------------------------------------------------------------------
% 2. MECCANISMI DI CONTROLLO
% -----------------------------------------------------------------------------
\section{Meccanismi di regolazione}
\label{sec:meccanismi}

I meccanismi di controllo si classificano in:

\begin{description}[leftmargin=1.5em,style=nextline]
  \item[Feedback (retroazione)] La variazione di un parametro dal valore fisiologico induce una reazione che agisce sulla causa stessa. Intervengono \textit{dopo} che un cambiamento è avvenuto.
  
  \item[Feedforward (anticipatorio)] Intervengono per controllare un processo \textit{prima} che avvenga una variazione (es. comando centrale nell'esercizio fisico).
  
  \item[Regolazioni comportamentali] Atteggiamenti volontari o appresi che contribuiscono al mantenimento dell'equilibrio (es. cercare ombra, bere acqua).
\end{description}

\begin{concetto}
\textit{«Il problema vero della fisiologia è la conoscenza dei meccanismi di regolazione dell'organismo»}.
\end{concetto}

\subsection{Il sistema a feedback}
\label{subsec:feedback}

Un sistema di controllo a feedback è costituito da tre elementi:
\begin{enumerate}
  \item \textbf{Sensori/Recettori}: rilevano la variabile controllata;
  \item \textbf{Centro di integrazione}: confronta il valore rilevato con il \textit{set point} e decide la risposta;
  \item \textbf{Effettori}: eseguono la risposta correttiva.
\end{enumerate}

I feedback si distinguono in:
\begin{itemize}
  \item \textbf{Feedback negativo}: la risposta contrasta la perturbazione iniziale (più comune, stabilizzante);
  \item \textbf{Feedback positivo}: la risposta amplifica la perturbazione (raro, es. parto, coagulazione).
\end{itemize}

\begin{figure}[htbp]
  \centering
  % Sostituire con diagramma reale se disponibile
  % \includegraphics[width=0.7\textwidth]{feedback_schema.pdf}
  \caption{Schema concettuale del controllo a feedback: rilevazione, integrazione, risposta.}
  \label{fig:feedback_schema}
\end{figure}

% -----------------------------------------------------------------------------
% 3. ESEMPIO APPLICATIVO: TERMOREGOLAZIONE
% -----------------------------------------------------------------------------
\section{Esempio: regolazione della temperatura corporea}
\label{sec:termoregolazione}

\begin{esempio}
\textbf{Mantenimento della temperatura interna a 37°C ± 0.5}
\end{esempio}

\begin{table}[htbp]
  \centering
  \caption{Risposte effettrici in base alla deviazione termica}
  \label{tab:termoregolazione}
  \begin{tabularx}{\textwidth}{l X}
    \toprule
    \textbf{Condizione} & \textbf{Risposte fisiologiche e comportamentali} \\
    \midrule
    Temperatura $> 37^\circ\text{C} \pm 0.5$ & 
    \begin{itemize}[leftmargin=1.5em,noitemsep,topsep=0pt,partopsep=0pt]
      \item Vasodilatazione periferica
      \item Attivazione ghiandole sudoripare
      \item Riduzione attività metabolica
      \item Comportamenti di dispersione (es. cercare fresco)
    \end{itemize} \\
    \addlinespace
    Temperatura $< 37^\circ\text{C} \pm 0.5$ & 
    \begin{itemize}[leftmargin=1.5em,noitemsep,topsep=0pt,partopsep=0pt]
      \item Vasocostrizione periferica
      \item Assenza di sudorazione
      \item Brivido e aumento metabolismo
      \item Comportamenti di conservazione (es. coprirsi, muoversi)
    \end{itemize} \\
    \bottomrule
  \end{tabularx}
\end{table}

\begin{remark}
Il centro di controllo è l'\textbf{ipotalamo}, che confronta i segnali dei termocettori centrali e periferici con il \textit{set point} e attiva le risposte appropriate.
\end{remark}

% -----------------------------------------------------------------------------
% 4. COMPARTIMENTI LIQUIDI E PARAMETRI OMEOSTATICI
% -----------------------------------------------------------------------------
\section{Ambiente interno e parametri critici}
\label{sec:ambiente_interno}

Attraverso l'ambiente extracellulare avvengono gli scambi di energia e sostanze con il mondo esterno. È necessario mantenere costanti parametri quali: pressione arteriosa, glicemia, pH, temperatura corporea, composizione e distribuzione dei liquidi ed elettroliti.

\subsection{Distribuzione dei liquidi corporei}
\label{subsec:liquidi}

\begin{table}[htbp]
  \centering
  \caption{Compartimenti liquidi in percentuale del peso corporeo}
  \label{tab:liquidi}
  \begin{tabularx}{0.85\textwidth}{l c c}
    \toprule
    \textbf{Compartimento} & \textbf{Sottocompartimento} & \textbf{\% peso corporeo} \\
    \midrule
    \multirow{2}{*}{Liquido extracellulare (LEC) \textasciitilde 20\%} 
      & Plasma sanguigno & 5\% \\
      & Liquido interstiziale & 15\% \\
    \addlinespace
    Liquido intracellulare (LIC) & Citoplasma & 40\% \\
    \bottomrule
  \end{tabularx}
\end{table}

% -----------------------------------------------------------------------------
% 5. ADATTAMENTO ALL'ESERCIZIO FISICO
% -----------------------------------------------------------------------------
\section{Adattamento e risposta all'esercizio}
\label{sec:adattamento}

\begin{concetto}
\textbf{Adattamento}: miglioramento delle condizioni fisiche e delle capacità funzionali. Piccoli aumenti di affaticamento e debolezza in seguito all'esercizio sono seguiti da \textit{ipercompensazione} durante il recupero.
\end{concetto}

Il ciclo adattativo segue la sequenza:
\[
\text{Esercizio} \rightarrow \text{Debilitazione transitoria} \rightarrow \text{Recupero} \rightarrow \text{Ipercompensazione} \rightarrow \text{Adattamento stabile}
\]

\subsection{Integrazione nervosa ed endocrina durante l'esercizio}
\label{subsec:integrazione}

\begin{table}[htbp]
  \centering
  \caption{Regolazione dei principali sistemi durante l'esercizio fisico}
  \label{tab:sistemi_esercizio}
  \begin{tabularx}{\textwidth}{l l l}
    \toprule
    \textbf{Sistema} & \textbf{Risposta} & \textbf{Variabile regolata} \\
    \midrule
    Circolatorio & Aumento frequenza e flusso & Pressione ematica \\
    Respiratorio & Aumento frequenza & Ossigenazione \\
    Rene & Riassorbimento acqua e sali & Volume ematico \\
    Fegato & Aumento immissione glucosio & Glicemia \\
    Muscolo & Contrazione-rilasciamento & -- \\
    Cute & Aumento sudorazione & Temperatura corporea \\
    Endocrino & Regolazione livelli ormonali & -- \\
    Digerente & Diminuzione attività & -- \\
    \bottomrule
  \end{tabularx}
\end{table}

\begin{esempio}
\textbf{Regolazione dell'osmolarità durante sudorazione intensa}:
\begin{enumerate}
  \item L'esercizio intenso provoca sudorazione $\rightarrow$ perdita di acqua plasmatica;
  \item Aumenta l'osmolarità del sangue $\rightarrow$ attivazione osmocettori ipotalamici;
  \item Risposte coordinate:
  \begin{itemize}
    \item Attivazione del centro della sete (stimolo all'ingestione di acqua);
    \item Secrezione di \textbf{ADH} (ormone antidiuretico) dalla neuroipofisi;
    \item L'ADH agisce sui tubuli renali $\rightarrow$ aumento riassorbimento di acqua;
    \item Risultato: riduzione dell'osmolarità ematica e ripristino dell'equilibrio idro-salino.
  \end{itemize}
\end{enumerate}
\end{esempio}

% -----------------------------------------------------------------------------
% 6. CONFRONTO: CONTROLLO NERVOSO VS ENDOCRINO
% -----------------------------------------------------------------------------
\section{Confronto tra sistemi di controllo}
\label{sec:confronto_controlli}

\begin{table}[htbp]
  \centering
  \caption{Differenze tra controllo nervoso ed endocrino}
  \label{tab:nervoso_vs_endocrino}
  \begin{tabularx}{\textwidth}{l X X}
    \toprule
    \textbf{Caratteristica} & \textbf{Sistema Nervoso} & \textbf{Sistema Endocrino} \\
    \midrule
    Trasmissione del segnale & Sinapsi (via neurale) & Sangue (via ematica) \\
    Tipo di segnale & Elettrochimico & Chimico (ormoni) \\
    Sede della risposta & Localizzata & Diffusa \\
    Velocità della risposta & Rapida (ms-s) & Lenta (s-min-ore) \\
    Durata della risposta & Breve & Duratura \\
    Processi controllati & Rapidi, precisi & Lenti, sostenuti \\
    \bottomrule
  \end{tabularx}
\end{table}

% =============================================================================
% CAPITOLO 2: ECCITABILITÀ CELLULARE
% Corso: Fisiologia dell'esercizio fisico
% =============================================================================

\chapter{Eccitabilità cellulare}
\label{chap:eccitabilita}

% -----------------------------------------------------------------------------
% 1. INTRODUZIONE: LA CELLULA COME UNITÀ FUNZIONALE
% -----------------------------------------------------------------------------
\section{La cellula eccitabile}
\label{sec:cellula_eccitabile}

\begin{definizione}
\textbf{Cellula eccitabile}: cellula (es. neurone, fibra muscolare) capace di reagire a uno stimolo adeguato con una variazione rapida e transitoria delle proprietà elettriche di membrana.
\end{definizione}

Ogni tipo cellulare è specializzato per compiti specifici, ma tutte le cellule condividono caratteristiche strutturali e funzionali di base. Le cellule eccitabili si distinguono per la capacità di generare e propagare segnali elettrici.

% -----------------------------------------------------------------------------
% 2. STRUTTURA E FUNZIONE DELLA MEMBRANA PLASMATICA
% -----------------------------------------------------------------------------
\section{La membrana plasmatica}
\label{sec:membrana}

\begin{concetto}
La membrana plasmatica è una struttura dinamica che:
\begin{itemize}
  \item è formata da un \textbf{bistrato lipidico} (isolante, idrofobico) con proteine e glicoproteine immerse;
  \item separa l'ambiente intracellulare da quello extracellulare;
  \item è \textbf{semipermeabile}: il passaggio di sostanze avviene secondo leggi chimico-fisiche (diffusione, trasporto facilitato, trasporto attivo).
\end{itemize}
\end{concetto}

\begin{figure}[htbp]
  \centering
  % \includegraphics[width=0.7\textwidth]{membrana_canali.pdf}
  \caption{Schema del mosaico fluido: canali ionici immersi nel bistrato fosfolipidico.}
  \label{fig:membrana_canali}
\end{figure}

\subsection{Distribuzione ionica e potenziale di equilibrio}
\label{subsec:distribuzione_ioni}

La differenza di concentrazione ionica ai due lati della membrana genera un potenziale elettrico. I principali ioni coinvolti sono:

\begin{table}[htbp]
  \centering
  \caption{Concentrazioni ioniche e potenziali di equilibrio (Nernst)}
  \label{tab:ioni_equilibrio}
  \begin{tabularx}{\textwidth}{l c c c}
    \toprule
    \textbf{Ione} & \textbf{[Intracellulare] (mM)} & \textbf{[Extracellulare] (mM)} & \textbf{Potenziale di equilibrio (mV)} \\
    \midrule
    \(\mathrm{K}^+\) & 140 & 5 & \(-75\) \\
    \(\mathrm{Na}^+\) & 5 & 140 & \(+55\) \\
    \(\mathrm{Cl}^-\) & 4.2 & 110 & \(-60\) \\
    \(\mathrm{A}^-\) (anioni organici) & \multicolumn{2}{c}{Non permeanti} & -- \\
    \bottomrule
  \end{tabularx}
\end{table}

\begin{remark}
Il potenziale di membrana a riposo (\(V_m\)) è circa \(-70\,\mathrm{mV}\), determinato principalmente dalla permeabilità al \(\mathrm{K}^+\) e dall'azione della pompa \(\mathrm{Na}^+\)/\(\mathrm{K}^+\).
\end{remark}

% -----------------------------------------------------------------------------
% 3. PROPRIETÀ ELETTRICHE DI MEMBRANA
% -----------------------------------------------------------------------------
\section{Proprietà passive e attive}
\label{sec:proprieta_elettriche}

\subsection{Risposte passive (elettrotoniche)}
\label{subsec:passive}

\begin{itemize}
  \item \textbf{Iperpolarizzazione}: variazione del potenziale verso valori più negativi di \(V_m\);
  \item \textbf{Depolarizzazione}: variazione verso valori meno negativi (più positivi) di \(V_m\);
  \item L'ampiezza della variazione è \textbf{proporzionale} all'intensità della corrente applicata.
\end{itemize}

Queste risposte sono \textit{graduate} e \textit{decrementali}: si attenuano con la distanza.

\subsection{Risposta attiva: il potenziale d'azione}
\label{subsec:potenziale_azione}

\begin{definizione}
\textbf{Potenziale d'azione}: variazione rapida, transitoria e \textit{tutto-o-nulla} del potenziale di membrana, innescata quando la depolarizzazione supera una soglia critica.
\end{definizione}

\begin{figure}[htbp]
  \centering
  % \includegraphics[width=0.8\textwidth]{potenziale_azione_fasi.pdf}
  \caption{Fasi del potenziale d'azione: (1-2) risposte passive graduate; (3) potenziale d'azione tutto-o-nulla.}
  \label{fig:pa_fasi}
\end{figure}

\subsubsection*{Fasi del potenziale d'azione}
\begin{enumerate}
  \item \textbf{Fase ascendente (depolarizzante)}: apertura rapida dei canali per il \(\mathrm{Na}^+\) voltaggio-dipendenti $\rightarrow$ ingresso massivo di \(\mathrm{Na}^+\) $\rightarrow$ inversione del potenziale;
  \item \textbf{Picco}: il potenziale raggiunge valori positivi (\(\sim +30\,\mathrm{mV}\));
  \item \textbf{Fase discendente (ripolarizzante)}: inattivazione dei canali \(\mathrm{Na}^+\) e apertura ritardata dei canali per il \(\mathrm{K}^+\) $\rightarrow$ uscita di \(\mathrm{K}^+\) $\rightarrow$ ritorno verso \(V_m\);
  \item \textbf{Iperpolarizzazione postuma}: temporanea eccessiva uscita di \(\mathrm{K}^+\) prima del ripristino dei gradienti.
\end{enumerate}

\begin{concetto}
Il potenziale d'azione è un fenomeno \textbf{tutto-o-nulla}: se lo stimolo supera la soglia, l'ampiezza e la forma sono indipendenti dall'intensità dello stimolo. L'informazione è codificata nella \textit{frequenza} degli impulsi, non nella loro ampiezza.
\end{concetto}

% -----------------------------------------------------------------------------
% 4. PERIODI REFRATTARI
% -----------------------------------------------------------------------------
\section{Periodi refrattari}
\label{sec:refrattari}

\begin{table}[htbp]
  \centering
  \caption{Caratteristiche dei periodi refrattari}
  \label{tab:refrattari}
  \begin{tabularx}{\textwidth}{l X}
    \toprule
    \textbf{Periodo} & \textbf{Caratteristiche funzionali} \\
    \midrule
    \textbf{Refrattario assoluto} & 
    \begin{itemize}[leftmargin=1.5em,noitemsep,topsep=0pt,partopsep=0pt]
      \item Impossibilità di generare un nuovo potenziale d'azione, indipendentemente dall'intensità dello stimolo;
      \item Corrisponde all'inattivazione dei canali \(\mathrm{Na}^+\);
      \item Garantisce la direzionalità unidirezionale dell'impulso.
    \end{itemize} \\
    \addlinespace
    \textbf{Refrattario relativo} & 
    \begin{itemize}[leftmargin=1.5em,noitemsep,topsep=0pt,partopsep=0pt]
      \item È possibile generare un potenziale d'azione solo con stimoli \textit{più intensi} del normale;
      \item Corrisponde alla fase di ripolarizzazione/iperpolarizzazione;
      \item Permette una modulazione della frequenza di scarica.
    \end{itemize} \\
    \bottomrule
  \end{tabularx}
\end{table}

% -----------------------------------------------------------------------------
% 5. PROPAGAZIONE DELL'IMPULSO NERVOSO
% -----------------------------------------------------------------------------
\section{Propagazione lungo l'assone}
\label{sec:propagazione}

\subsection{Conduzione continua (assoni non mielinizzati)}
\label{subsec:continua}
L'impulso si rigenera punto per punto lungo tutta la membrana assonica. È un processo \textbf{lento} (0.5–2 m/s) e metabolicamente costoso.

\subsection{Conduzione saltatoria (assoni mielinizzati)}
\label{subsec:saltatoria}

\begin{concetto}
La \textbf{guaina mielinica}, formata da cellule della glia (Schwann nel SNP, oligodendrociti nel SNC), isola elettricamente l'assone. L'impulso "salta" da un \textbf{nodo di Ranvier} al successivo, dove sono concentrati i canali ionici voltaggio-dipendenti.
\end{concetto}

\begin{itemize}
  \item \textbf{Vantaggi}: velocità fino a \(120\,\mathrm{m/s}\), minore dispersione di corrente, risparmio energetico;
  \item \textbf{Implicazioni cliniche}: la demielinizzazione (es. sclerosi multipla) interrompe la conduzione saltatoria, compromettendo la trasmissione nervosa.
\end{itemize}

\begin{figure}[htbp]
  \centering
  % \includegraphics[width=0.75\textwidth]{conduzione_saltatoria.pdf}
  \caption{Conduzione saltatoria: l'impulso si rigenera solo ai nodi di Ranvier (A: innesco; B: progressione).}
  \label{fig:saltatoria}
\end{figure}

% -----------------------------------------------------------------------------
% 6. TRASMISSIONE SINAPTICA (INTRODUZIONE)
% -----------------------------------------------------------------------------
\section{La sinapsi: trasmissione dell'informazione}
\label{sec:sinapsi_intro}

\begin{definizione}
\textbf{Sinapsi}: struttura specializzata che permette la trasmissione dell'informazione da un neurone (o cellula eccitabile) a un'altra cellula bersaglio.
\end{definizione}

Il potenziale d'azione, una volta generato, deve essere trasmesso ad altre cellule per coordinare funzioni fisiologiche. Questo passaggio avviene attraverso:
\begin{itemize}
  \item \textbf{Sinapsi elettriche}: trasmissione diretta tramite giunzioni serrate (gap junctions);
  \item \textbf{Sinapsi chimiche}: rilascio di neurotrasmettitori che legano recettori postsinaptici (più comuni e modulabili).
\end{itemize}

\begin{esempio}
\textbf{Esempio applicativo – Contrazione muscolare}:
\begin{enumerate}
  \item Un potenziale d'azione viaggia lungo il motoneurone $\alpha$;
  \item Raggiunge la terminazione sinaptica sulla placca motrice;
  \item Il neurotrasmettitore (acetilcolina) viene rilasciato e lega i recettori nicotinici sulla fibra muscolare;
  \item Si genera un potenziale di placca che, se supera la soglia, innesca un potenziale d'azione muscolare $\rightarrow$ contrazione.
\end{enumerate}
\end{esempio}

% =============================================================================
% CAPITOLO: TRASMISSIONE SINAPTICA
% Corso: Fisiologia dell'esercizio fisico
% =============================================================================

\chapter{Trasmissione sinaptica}
\label{chap:trasmissione_sinaptica}

% -----------------------------------------------------------------------------
% 1. MODALITÀ DI COMUNICAZIONE CELLULARE
% -----------------------------------------------------------------------------
\section{Modalità di comunicazione cellulare}
\label{sec:comunicazione_cellulare}

Le cellule coordinano le loro attività attraverso quattro principali modalità di segnalazione:

\begin{description}[leftmargin=1.5em,style=nextline]
  \item[Nervosa] Segnale elettrico che comporta il rilascio di un neurotrasmettitore sulla cellula bersaglio. Tipica del sistema nervoso.
  \item[Endocrina] Segnale chimico: ormoni secreti nel sangue raggiungono cellule bersaglio lontane, dotate di recettori specifici. L'interazione ormone-recettore attiva risposte come sintesi/attivazione di proteine o modulazione di canali ionici.
  \item[Neuroendocrina] Segnale elettrico che induce il rilascio di ormoni nel circolo sanguigno.
  \item[Paracrina e autocrina] Segnale chimico rilasciato nel liquido extracellulare (LEC) a breve distanza, agendo su cellule vicine (paracrina) o sulla stessa cellula secernente (autocrina).
\end{description}

\begin{concetto}
La comunicazione cellulare è il fondamento dell'integrazione fisiologica. Il sistema nervoso privilegia velocità e precisione spaziale, mentre quello endocrino garantisce ampiezza e durata d'azione.
\end{concetto}

% -----------------------------------------------------------------------------
% 2. LE SINAPSI: ARCHITETTURA E TIPOLOGIE
% -----------------------------------------------------------------------------
\section{Le sinapsi: architettura e tipologie}
\label{sec:sinapsi}

\begin{definizione}
\textbf{Sinapsi}: punto di contatto \textit{funzionale} (non necessariamente fisico) tra due cellule, specializzato nella trasmissione del segnale da un elemento presinaptico a uno postsinaptico.
\end{definizione}

\subsection{Sinapsi elettriche}
\label{subsec:sinapsi_elettriche}

Costituite dall'accostamento delle membrane di due cellule a una distanza di 3–4 nm, formano una \textbf{giunzione comunicante} (\textit{gap junction}). Questa connessione a bassa resistenza permette il passaggio diretto e passivo di ioni.

\begin{itemize}
  \item \textbf{Velocità}: trasmissione senza ritardo;
  \item \textbf{Direzionalità}: bidirezionale, dipende dal gradiente ionico;
  \item \textbf{Segnale}: induce nel postsinaptico variazioni dello stesso segno del presinaptico;
  \item \textbf{Localizzazione}: nuclei ipotalamici neuroendocrini, miociti cardiaci/lisci.
\end{itemize}

\subsection{Sinapsi chimiche}
\label{subsec:sinapsi_chimiche}

Presentano una struttura più complessa, organizzata in tre componenti:
\begin{enumerate}
  \item \textbf{Elemento presinaptico}: contiene mitocondri, citoscheletro e vescicole sinaptiche cariche di neurotrasmettitore;
  \item \textbf{Fessura sinaptica}: spazio di 20–40 nm;
  \item \textbf{Elemento postsinaptico}: membrana con recettori specifici per il neurotrasmettitore.
\end{enumerate}

Possono connettere neurone-neurone, neurone-muscolo o neurone-ghiandola, e classificarsi morfologicamente in \textit{asso-dendritiche}, \textit{asso-somatiche} o \textit{asso-assoniche}.

\begin{concetto}
A differenza delle sinapsi elettriche, la trasmissione chimica presenta:
\begin{itemize}[leftmargin=1.2em,noitemsep,topsep=0pt,partopsep=0pt]
  \item \textbf{Ritardo sinaptico}: $\sim 0.5\,\mathrm{ms}$ per rilascio, diffusione e legame;
  \item \textbf{Unidirezionalità}: dal presinaptico al postsinaptico;
  \item \textbf{Plasticità}: può generare potenziali eccitatori (EPSP) o inibitori (IPSP).
\end{itemize}
\end{concetto}

% -----------------------------------------------------------------------------
% 3. INTEGRAZIONE POSTSINAPTICA E POTENZIALI GRADUATI
% -----------------------------------------------------------------------------
\section{Integrazione postsinaptica e potenziali graduati}
\label{sec:integrazione_posts}

Il legame neurotrasmettitore-recettore modifica la permeabilità di membrana, generando potenziali graduati:
\begin{itemize}
  \item \textbf{EPSP} (\textit{Excitatory Post-Synaptic Potential}): depolarizzazione locale che avvicina la membrana alla soglia del potenziale d'azione;
  \item \textbf{IPSP} (\textit{Inhibitory Post-Synaptic Potential}): iperpolarizzazione o stabilizzazione che allontana la membrana dalla soglia.
\end{itemize}

\subsection{Sommazione spaziale e temporale}
\label{subsec:sommazione}

Un singolo EPSP raramente supera la soglia. L'integrazione avviene tramite:
\begin{description}[leftmargin=1.5em,style=nextline]
  \item[Sommazione spaziale] Attivazione simultanea di più sinapsi eccitatorie su diversi dendriti. Gli EPSP si sommano algebricamente; la presenza contemporanea di IPSP può contrastare l'effetto.
  \item[Sommazione temporale] Attivazione ripetuta e ravvicinata della stessa sinapsi. Gli EPSP successivi si sommano prima che il precedente si sia esaurito.
\end{description}

\begin{esempio}
\textbf{Dal potenziale graduato al potenziale d'azione}:
Se la somma algebrica di EPSP e IPSP a livello del cono di emergenza (assone iniziale) supera il valore di soglia ($\sim -55\,\mathrm{mV}$), si innesca un potenziale d'azione che viaggia lungo l'assone. In caso contrario, il segnale si estingue.
\end{esempio}

% -----------------------------------------------------------------------------
% 4. NEUROTRASMETTITORI E RECETTORI
% -----------------------------------------------------------------------------
\section{Neurotrasmettitori e recettori}
\label{sec:neurotrasmettitori}

\subsection{Classificazione dei neurotrasmettitori}
\label{subsec:nt_classificazione}

I neurotrasmettitori (NT) si suddividono per struttura chimica:

\begin{table}[htbp]
  \centering
  \caption{Principali classi di neurotrasmettitori}
  \label{tab:nt_classi}
  \begin{tabularx}{\textwidth}{l X}
    \toprule
    \textbf{Classe} & \textbf{Esempi e localizzazione} \\
    \midrule
    Molecole a basso peso molecolare & Acetilcolina (ACh), amine biogene (dopamina, serotonina, noradrenalina) \\
    \addlinespace
    Aminoacidi & Glutammato (eccitatorio), GABA e glicina (inibitori), aspartato \\
    \addlinespace
    Peptidi & Oltre 100 tipi (es. endorfine, sostanza P), catene di 3–30 aminoacidi \\
    \addlinespace
    Gas e purine & ATP, adenosina, ossido nitrico (NO), monossido di carbonio (CO) \\
    \bottomrule
  \end{tabularx}
\end{table}

\begin{remark}
Nel sistema nervoso periferico somatico, il NT principale è l'\textbf{acetilcolina} (giunzione neuromuscolare). Nel SNC operano reti complesse con glutammato, GABA, dopamina e serotonina.
\end{remark}

\subsection{Recettori postsinaptici}
\label{subsec:recettori}

L'effetto del NT dipende dal tipo di recettore:
\begin{itemize}
  \item \textbf{Ionotropi}: canale ionico intrinseco. Legame rapido $\rightarrow$ apertura diretta $\rightarrow$ risposta immediata (EPSP/IPSP);
  \item \textbf{Metabotropi}: accoppiati a proteine G. Attivano cascate di secondi messaggeri (AMP ciclico, $\mathrm{Ca}^{2+}$, chinasi) $\rightarrow$ effetti più lenti, duraturi e modulanti (es. regolazione genica, plasticità sinaptica).
\end{itemize}

% -----------------------------------------------------------------------------
% 5. RUOLO FUNZIONALE DELLA SINAPSI CHIMICA
% -----------------------------------------------------------------------------
\section{Ruolo funzionale della sinapsi chimica}
\label{sec:ruolo_funzionale}

Le sinapsi chimiche non sono semplici ``interruttori'', ma veri e propri \textbf{modulatori del segnale}. Grazie alla loro architettura, consentono:
\begin{enumerate}
  \item \textbf{Amplificazione}: un singolo NT può attivare molte proteine G o aprire numerosi canali;
  \item \textbf{Inversione}: un segnale eccitatorio può diventare inibitorio a seconda del recettore espresso;
  \item \textbf{Prolungamento temporale}: le cascate metabotropiche mantengono la risposta ben oltre la presenza del NT;
  \item \textbf{Plasticità sinaptica}: modificazioni strutturali e funzionali a lungo termine (base dell'apprendimento e della memoria).
\end{enumerate}


% =============================================================================
% CAPITOLO: ORGANIZZAZIONE DEL SISTEMA NERVOSO
% Corso: Fisiologia dell'esercizio fisico
% =============================================================================

\chapter{Organizzazione del sistema nervoso}
\label{chap:org_sistema_nervoso}

% -----------------------------------------------------------------------------
% 1. FUNZIONE INTEGRATIVA DEL SISTEMA NERVOSO
% -----------------------------------------------------------------------------
\section{La funzione integrativa del sistema nervoso}
\label{sec:funzione_integrativa}

\begin{definizione}
\textbf{Funzione integrativa}: capacità del sistema nervoso di elaborare informazioni sensoriali, coordinare risposte motorie e regolative, e mantenere l'omeostasi dell'organismo attraverso circuiti neurali specializzati.
\end{definizione}

Il sistema nervoso rappresenta il principale sistema di controllo dell'organismo, caratterizzato da:
\begin{itemize}[leftmargin=1.2em,noitemsep,topsep=0pt,partopsep=0pt]
  \item \textbf{Rapidità}: trasmissione elettrica in millisecondi;
  \item \textbf{Precisione spaziale}: connessioni sinaptiche specifiche;
  \item \textbf{Plasticità}: capacità di modificare l'efficacia delle connessioni in base all'esperienza.
\end{itemize}

\begin{concetto}
Il sistema nervoso non opera in isolamento: l'integrazione con il sistema endocrino garantisce risposte fisiologiche coordinate e adattative, fondamentali durante l'esercizio fisico.
\end{concetto}

% -----------------------------------------------------------------------------
% 2. ORGANIZZAZIONE ANATOMO-FUNZIONALE
% -----------------------------------------------------------------------------
\section{Organizzazione anatomo-funzionale}
\label{sec:organizzazione_anatomo_funzionale}

Il sistema nervoso si articola in due grandi compartimenti:

\begin{table}[htbp]
  \centering
  \caption{Suddivisione del sistema nervoso}
  \label{tab:divisione_sn}
  \begin{tabularx}{\textwidth}{l X}
    \toprule
    \textbf{Compartimento} & \textbf{Componenti e funzioni} \\
    \midrule
    \textbf{Sistema Nervoso Centrale (SNC)} & 
    \begin{itemize}[leftmargin=1.2em,noitemsep,topsep=0pt,partopsep=0pt]
      \item \textit{Encefalo}: corteccia cerebrale, sistema limbico, ipotalamo, tronco encefalico, cervelletto;
      \item \textit{Midollo spinale}: via di transito e centro di integrazione riflessa;
      \item Funzione: elaborazione centrale, coordinamento, memoria, controllo volontario.
    \end{itemize} \\
    \addlinespace
    \textbf{Sistema Nervoso Periferico (SNP)} & 
    \begin{itemize}[leftmargin=1.2em,noitemsep,topsep=0pt,partopsep=0pt]
      \item \textit{Nervi cranici e spinali}: conduzione afferente (sensoriale) ed efferente (motoria);
      \item \textit{Gangli}: aggregazioni di corpi neuronali esterni al SNC;
      \item Funzione: collegamento tra SNC e organi/tessuti periferici.
    \end{itemize} \\
    \bottomrule
  \end{tabularx}
\end{table}

Il SNP si suddivide funzionalmente in:
\begin{description}[leftmargin=1.5em,style=nextline]
  \item[Sistema somatico] Controllo volontario della muscolatura scheletrica e ricezione di stimoli sensoriali esterni.
  \item[Sistema autonomo (viscerale)] Regolazione involontaria di organi interni, ghiandole e muscolatura liscia/cardica. Si articola in:
  \begin{itemize}[leftmargin=1.2em,noitemsep,topsep=0pt,partopsep=0pt]
    \item \textbf{Simpatico} (toraco-lombare): prepara l'organismo all'azione (``lotta o fuga'');
    \item \textbf{Parasimpatico} (cranio-sacrale): promuove il risparmio energetico e le funzioni di ``riposo e digestione'';
    \item \textbf{Enterico}: rete intrinseca del tratto gastrointestinale, parzialmente autonoma.
  \end{itemize}
\end{description}

% -----------------------------------------------------------------------------
% 3. LE CELLULE DEL SISTEMA NERVOSO: NEURONI E GLIA
% -----------------------------------------------------------------------------
\section{Cellule del sistema nervoso}
\label{sec:cellule_sn}

\subsection{I neuroni: unità funzionali}
\label{subsec:neuroni}

\begin{definizione}
\textbf{Neurone}: cellula eccitabile specializzata nella ricezione, integrazione e trasmissione di segnali elettrochimici.
\end{definizione}

Struttura tipica di un neurone multipolare:
\begin{itemize}
  \item \textbf{Corpo cellulare (soma)}: contiene il nucleo e gli organelli metabolici;
  \item \textbf{Dendriti}: prolungamenti riceventi, aumentano la superficie sinaptica;
  \item \textbf{Assone}: prolungamento conduttore, trasmette il potenziale d'azione verso le terminazioni;
  \item \textbf{Terminazioni sinaptiche}: rilasciano neurotrasmettitori sulla cellula bersaglio.
\end{itemize}

\begin{figure}[htbp]
  \centering
  % \includegraphics[width=0.7\textwidth]{neurone_struttura.pdf}
  \caption{Schema di un neurone multipolare: dendriti, soma, assone mielinizzato, terminazioni.}
  \label{fig:neurone_struttura}
\end{figure}

\subsection{Le cellule gliali: supporto e modulazione}
\label{subsec:glia}

\begin{concetto}
Le \textbf{cellule gliali} (o \textit{nevroglia}) superano numericamente i neuroni (rapporto $\sim 10:1$) e svolgono funzioni essenziali di supporto strutturale, metabolico e funzionale.
\end{concetto}

\begin{table}[htbp]
  \centering
  \caption{Principali tipi di cellule gliali e funzioni}
  \label{tab:cellule_gliali}
  \begin{tabularx}{\textwidth}{l X}
    \toprule
    \textbf{Tipo cellulare} & \textbf{Localizzazione e funzioni principali} \\
    \midrule
    \textbf{Astrociti} & 
    \begin{itemize}[leftmargin=1.2em,noitemsep,topsep=0pt,partopsep=0pt]
      \item SNC; supporto strutturale, regolazione del microambiente ionico, barriera emato-encefalica, riserva energetica (glicogeno);
    \end{itemize} \\
    \addlinespace
    \textbf{Oligodendrociti} & 
    \begin{itemize}[leftmargin=1.2em,noitemsep,topsep=0pt,partopsep=0pt]
      \item SNC; formano la guaina mielinica attorno a più assoni, accelerando la conduzione;
    \end{itemize} \\
    \addlinespace
    \textbf{Cellule di Schwann} & 
    \begin{itemize}[leftmargin=1.2em,noitemsep,topsep=0pt,partopsep=0pt]
      \item SNP; mielinizzazione degli assoni periferici, supporto alla rigenerazione nervosa;
    \end{itemize} \\
    \addlinespace
    \textbf{Microglia} & 
    \begin{itemize}[leftmargin=1.2em,noitemsep,topsep=0pt,partopsep=0pt]
      \item SNC; funzione immunitaria, fagocitosi di detriti, risposta infiammatoria;
    \end{itemize} \\
    \addlinespace
    \textbf{Cellule ependimali} & 
    \begin{itemize}[leftmargin=1.2em,noitemsep,topsep=0pt,partopsep=0pt]
      \item Rivestimento dei ventricoli cerebrali e del canale ependimale; produzione e circolazione del liquido cerebrospinale.
    \end{itemize} \\
    \bottomrule
  \end{tabularx}
\end{table}

\begin{remark}
La \textbf{guaina mielinica}, formata da oligodendrociti (SNC) o cellule di Schwann (SNP), isola elettricamente l'assone e permette la \textit{conduzione saltatoria}, aumentando la velocità di propagazione fino a $120\,\mathrm{m/s}$ e riducendo il costo energetico.
\end{remark}

% -----------------------------------------------------------------------------
% 4. CIRCUITI NEURONALI E INTEGRAZIONE FUNZIONALE
% -----------------------------------------------------------------------------
\section{Circuiti neuronali}
\label{sec:circuiti_neuronali}

I neuroni non operano isolatamente, ma sono organizzati in \textbf{circuiti} che elaborano l'informazione attraverso schemi di connessione specifici:

\begin{description}[leftmargin=1.5em,style=nextline]
  \item[Circuito a convergenza] Più neuroni presinaptici convergono su un singolo neurone postsinaptico. Permette l'integrazione di segnali multipli (es. sommazione spaziale).
  \item[Circuito a divergenza] Un neurone presinaptico si connette a molti neuroni postsinaptici. Amplifica e distribuisce il segnale a più vie.
  \item[Circuito a feedback (ricorrente)] L'output di un neurone ritorna, direttamente o tramite interneuroni, a modulare la propria attività. Base dei meccanismi di regolazione e oscillazione.
  \item[Circuito a inibizione laterale] Un neurone eccitato inibisce i neuroni adiacenti, aumentando il contrasto e la precisione spaziale del segnale (es. elaborazione sensoriale).
\end{description}

\begin{esempio}
\textbf{Riflesso miotatico (da stiramento)}:
\begin{enumerate}
  \item Stiramento del muscolo $\rightarrow$ attivazione dei fusi neuromuscolari;
  \item Segnale afferente al midollo spinale $\rightarrow$ sinapsi eccitatoria diretta sul motoneurone $\alpha$ dello stesso muscolo (contrazione);
  \item Contemporaneamente, attivazione di un interneurone inibitorio sul motoneurone del muscolo antagonista (rilassamento);
  \item Risultato: correzione rapida della lunghezza muscolare, fondamentale per postura e coordinazione durante il movimento.
\end{enumerate}
\end{esempio}

% -----------------------------------------------------------------------------
% 5. IMPLICAZIONI PER LA FISIOLOGIA DELL'ESERCIZIO
% -----------------------------------------------------------------------------
\section{Integrazione nervosa durante l'esercizio}
\label{sec:esercizio_sn}

\begin{concetto}
Durante l'attività fisica, il sistema nervoso coordina in tempo reale funzioni multiple: attivazione motoria, regolazione cardiovascolare e respiratoria, percezione sensoriale, termoregolazione e modulazione endocrina.
\end{concetto}

\begin{table}[htbp]
  \centering
  \caption{Principali adattamenti neurali all'esercizio fisico}
  \label{tab:adattamenti_neurali}
  \begin{tabularx}{\textwidth}{l X}
    \toprule
    \textbf{Livello} & \textbf{Adattamento funzionale} \\
    \midrule
    \textbf{Corticale} & 
    \begin{itemize}[leftmargin=1.2em,noitemsep,topsep=0pt,partopsep=0pt]
      \item Miglioramento del reclutamento motorio e della coordinazione intermuscolare;
      \item Plasticità sinaptica nelle aree motorie (apprendimento del gesto sportivo).
    \end{itemize} \\
    \addlinespace
    \textbf{Spinale} & 
    \begin{itemize}[leftmargin=1.2em,noitemsep,topsep=0pt,partopsep=0pt]
      \item Modulazione dei riflessi (es. facilitazione del riflesso da stiramento);
      \item Riduzione dell'inibizione presinaptica per ottimizzare l'output motorio.
    \end{itemize} \\
    \addlinespace
    \textbf{Autonomo} & 
    \begin{itemize}[leftmargin=1.2em,noitemsep,topsep=0pt,partopsep=0pt]
      \item Bilanciamento simpatico/parasimpatico per adattare frequenza cardiaca, pressione e flusso ematico;
      \item Regolazione della sudorazione e del flusso cutaneo per la termoregolazione.
    \end{itemize} \\
    \bottomrule
  \end{tabularx}
\end{table}


\end{document}